\subsection{Capacitor}
A capacitor is a two-port passive device that stores energy in the form of
an electric field. The property of a capacitor is called capacitance and is
measured in Farads (F). The core relationship between charge, capacitance
and voltage is given by:
\begin{align*}
  Q = C \cdot V \\
  \frac{dQ}{dt} &= C \cdot \frac{dV}{dt} \hspace{1cm} \text{(assuming \(C\) is constant)} \\
  I &= C \cdot \frac{dV}{dt} \hspace{1cm} [I = \frac{dQ}{dt}]
\end{align*}
This is the fundamental relationship between current and voltage in a
capacitor or if you may call this the capacitor's version of Ohm's law.
The inverse of capacitance is called elastance (\(S\)) and is measured in Farads
inverse (F\(^{-1}\)).

The impedance of a capacitor is given by:
\[
  Z_C = \frac{1}{j\omega C} = \frac{1}{j 2 \pi f C} = \frac{1}{s C}
\]
The inverse of impedance is called admittance (\(Y\)) and is measured in
Siemens (S). The relationship between impedance and admittance is:
\[
  Y = \frac{1}{Z} = \frac{1}{\frac{1}{sC}} = sC
\]
You'll notice when we apply KCL to circuits within this book, we'll mostly
use admittance \(Y\) to express the impedance, we do this as it simplifies
the analysis a lot. The admittance version of Ohm's law is:
\[
  I = V \cdot Y
\]
The impedance of a capacitor is inversely proportional to frequency. This
means that at low frequencies, the impedance of a capacitor is very high and
at frequencies, the impedance of a capacitor is very low. This is why
capacitor effects are more pronounced at high frequencies as their
impedance drops enough to load the circuit. This is also why capacitors are
used for decoupling (they act as a short circuit at high frequencies).

An ideal capacitor is a linear (that is, the relationship between voltage and
current is a straight line) and time-invariant (that is, the relationship
between voltage and current does not change with time) of the three passive
components. It is also the only device that in its ideal form does not
dissipate energy.

A capacitor can be formed anywhere two conductors are separated by an
insulator. The charges are stored on the surface of the  conductors and the
energy is stored in the electric field between the  conductors. The energy
stored in a capacitor is given by:
\begin{align*}
  I &= C \cdot \frac{dV}{dt} \\
  P &= V \cdot I \\
  P &= V \cdot C \cdot \frac{dV}{dt} \\
  P &= C \cdot V \cdot \frac{dV}{dt} \\
  P &= C \cdot \frac{1}{2} \cdot \frac{d(V^2)}{dt} \hspace{1cm} [\frac{d(V^2)}{dt} = 2V \cdot \frac{dV}{dt}] \\
  \int P dt &= \int C \cdot \frac{1}{2} \cdot \frac{d(V^2)}{dt} dt \\
  E &= \frac{1}{2} C V^2
\end{align*}
